\begin{textsample}{POS Dim 1 – vem\_ed\_11 – Score 12.00 – vem\_ed\_11/t049.txt}  \label{ex:f1_pos_148}
BOAS PRÁTICAS Pedro Sérgio Rosa \textbf{é} doutor em Educação para a Ciência (Unesp), mestre em Física (Unicamp), Bacharel (UEL) e Licenciado (Unesp) em Física.
Professor do curso de Produção Fonográfica na Fatec Tatuí, relata como superar a barreira do \textbf{idioma} e desenvolver um Projeto Colaborativo Internacional em uma área aparentemente árida, a física das ondas.
Os desafios \textbf{são} vários, mas, o principal \textbf{é}, em tese, a barreira do \textbf{idioma}.
Por \textbf{que} em tese?
Porque \textbf{são} três aspectos: leitura, escrita e comunicação oral.
A leitura para mim \textbf{era} normal, pois, desde o Bacharelado em Física estava acostumado com os textos científicos em inglês.
Por outro lado, a escrita em outro \textbf{idioma} não \textbf{é} nada trivial, uma vez \textbf{que} \textbf{é} necessário estar continuamente estudando.
A \textbf{parte} mais complexa \textbf{é} a comunicação oral: compreender o \textbf{que} um nativo da língua fala e responder adequadamente.
Classifico por \textbf{nível} de dificuldade, a \textbf{partir} do mais fácil: leitura, escrita e comunicação oral em língua inglesa.
Para um nativo, a comunicação entre pares \textbf{é} sempre muito mais rápida.
Em relação ao ensino de Acústica Aplicada à Produção Fonográfica, \textbf{disciplina} \textbf{que} ministro na Fatec, observo \textbf{que}, para planejar um Projeto Colaborativo Internacional (PCI), \textbf{são} necessárias muitas reuniões, \textbf{diálogos} e correções de \textbf{rota}.
O \textbf{que} me ajudou enormemente \textbf{foi} ter realizado, com a professora de inglês Dulce Helena Soares Villa Nova e com o estudante Flavio Lee, da monitoria em línguas da Fatec Tatuí, aulas e correções na pronúncia, bem como treinos \textbf{que} \textbf{tornaram} minha audição mais aguda.
Isso \textbf{permitiu}, com o tempo, melhorar minha \textbf{compreensão} da fala de Machele Kindle e James Goodwin, professores colaboradores no PCI desenvolvido entre Fatec Tatuí e BridgeValley Community College (EUA).
Porém, como \textbf{todo} \textbf{processo} de aprimoramento, para se ter proficiência em outra língua, a quantidade de horas de prática \textbf{é} decisiva.
O contato com o nativo \textbf{cria} uma memória de audibilidade \textbf{que} \textbf{permite} assistir a um \textbf{programa}, filme, documentário na língua inglesa sem a \textbf{necessidade} da leitura total da tradução pela legenda.
Não há \textbf{que} se ter receio de \textbf{aprender} algo novo ou de adquirir novas \textbf{habilidades}.
A dica \textbf{é} não ter medo de errar, mas, se ocorrer o erro, pedir \textbf{que} o professor nativo ou o professor de línguas \textbf{que} participa do PCI ajude.
Durante o projeto, passei a assistir a filmes com áudio e legenda em inglês, para treinar a audição: creio \textbf{que} isso ajudou também.
Finalizando, diria \textbf{que} os Projetos Colaborativos Internacionais da Cesu / Centro Paula Souza precisam \textbf{ser} explorados pela comunidade de alunos e professores \textbf{que} consideram a cooperação como um ganho fundamental para seu aprimoramento, tanto pessoal \textbf{quanto} \textbf{profissional}.

% matched lemmas: aprender, compreensão, criar, disciplina, diálogo, habilidade, idioma, necessidade, nível, parte, partir, permitir, processo, profissional, programa, quanto, que, rota, ser, todo, tornar
\end{textsample}
